\documentclass[14pt]{extarticle}
% ---------- PACKAGES ----------
\usepackage{amsmath, amssymb}
\usepackage{geometry}
\usepackage{setspace}
\usepackage{titlesec}
\usepackage{xcolor}
\usepackage{helvet}
\usepackage{enumitem}
\usepackage{lmodern}
\makeatletter
\IfFileExists{../common-things/reflectionpage.tex}{%
  \input{../common-things/reflectionpage.tex}%
}{%
  \IfFileExists{common-things/reflectionpage.tex}{%
    \input{common-things/reflectionpage.tex}%
  }{%
    \PackageError{\jobname}{Missing reflectionpage.tex file}{%
      Place reflectionpage.tex inside a common-things/ directory next to this unit\@.%
    }%
  }%
}
\makeatother
% ---------- PAGE SETUP ----------
\geometry{margin=1in}
\setstretch{1.5}
\setlength{\parindent}{0pt}
\setlength{\parskip}{0.75em}
\setlist{itemsep=0.75\baselineskip, topsep=0.5\baselineskip}
\titleformat{\section}{\normalfont\Large\bfseries}{\thesection}{1em}{}
\titleformat{\subsection}{\normalfont\large\bfseries}{\thesubsection}{1em}{}
\renewcommand{\familydefault}{\sfdefault}
\renewcommand{\emph}[1]{\textbf{#1}}

\begin{document}
\raggedright
\pagenumbering{gobble}

\begin{center}
    \LARGE \textbf{Unit 10: Review and Test Strategy} \\[6pt]
    \Large \textbf{Topic 6: Full Mixed Review}
\end{center}

\vspace{1em}

\section*{Concept Summary}

This final topic is a comprehensive review that pulls from every unit in the curriculum. It simulates the range and variety of a full SAT Math section. There is no new content here. Instead, the 25 practice questions cover algebra, functions, geometry, trigonometry, data analysis, and probability in unpredictable order, just as they would appear on the real test.

Use this topic as a diagnostic: after working through the questions, note which topics caused you difficulty and revisit the corresponding unit for targeted review. The questions are arranged in approximate order of increasing difficulty, matching the structure of an SAT module.

\section*{Core Skills}
\begin{itemize}
  \item Apply the appropriate concept or formula without topic labels as prompts.
  \item Manage time effectively across a varied question set.
  \item Recognize connections between different math domains within a single problem.
  \item Verify answers by substitution, estimation, or graphing when time permits.
\end{itemize}

\section*{Key Takeaways}
\begin{itemize}
  \item This topic has no new formulas or techniques. Every question draws on skills from Units 1 through 9.
  \item Use your results here as a diagnostic. A missed question points back to a specific unit worth revisiting, not a gap in this topic itself.
  \item On the real SAT, questions do not announce their topic. Practicing without labels, as in this set, builds the recognition skill that matters most on test day.
  \item If you consistently run out of time before finishing, revisit Topic 4 (Pacing and the Two-Pass Approach) rather than just doing more practice questions.
\end{itemize}

\newpage

\section*{Practice Questions: Full Mixed Review}

\subsection*{Part A: Fundamentals Review}
\begin{enumerate}
  \item Solve for \(x\): \(4(x - 3) = 2x + 10\).
  \item If \(f(x) = x^2 - 3x + 2\), what is \(f(4)\)?
  \item A line passes through the points \((1, 3)\) and \((4, 12)\). What is the slope of the line?
  \item What is \(20\%\) of 150?
  \item The angles of a triangle measure \(x^\circ\), \(2x^\circ\), and \(60^\circ\). What is the value of \(x\)?
\end{enumerate}

\subsection*{Part B: Core Applications}
\begin{enumerate}
  \item Solve the system: \(3x + 2y = 16\) and \(x - y = 2\).
  \item A circle has a radius of 8. What is the area of a sector with a central angle of \(45^\circ\)? Express your answer in terms of \(\pi\).
  \item Simplify \(\dfrac{2x^2 - 8}{2x + 4}\) for \(x \neq -2\).
  \item The probability of drawing a red marble from a jar is \(\dfrac{3}{10}\). If there are 40 marbles, how many are red?
  \item A right triangle has legs of length 7 and 24. What is the length of the hypotenuse?
\end{enumerate}

\subsection*{Part C: Multi-Step Problems}
\begin{enumerate}
  \item A store sells an item for \(\$60\) after a \(25\%\) markup. What was the wholesale cost?
  \item The vertex of the parabola \(y = 2x^2 - 12x + 22\) is at \((h, k)\). What is the value of \(h + k\)?
  \item In a 30-60-90 triangle, the hypotenuse is 20. What is the length of the side opposite the \(60^\circ\) angle? Express your answer in simplified radical form.
  \item The mean of a data set of 8 values is 15. A ninth value of 33 is added. What is the new mean?
  \item Two similar triangles have a scale factor of \(\dfrac{2}{3}\). If the area of the larger triangle is 54, what is the area of the smaller triangle?
\end{enumerate}

\subsection*{Part D: Challenging Applications}
\begin{enumerate}
  \item If \(\sin(\theta) = \dfrac{8}{17}\) and \(\theta\) is acute, what is \(\tan(\theta)\)?
  \item A cylindrical container has a radius of 4 cm and is filled with water to a height of 9 cm. All of the water is poured into a second cylindrical container with a radius of 6 cm. What is the height of the water in the second container, in cm?
  \item The line \(y = -x + c\) is tangent to \(y = x^2 + 4x + 5\). What is the value of \(c\)?
  \item A population of 5{,}000 grows by \(6\%\) per year. Write an expression for the population after \(t\) years, and find the population after 5 years, to the nearest whole number.
  \item A rectangle is inscribed in a circle of radius 5. If one side of the rectangle has length 6, what is the area of the rectangle?
\end{enumerate}

\subsection*{Part E: SAT-Hard}
\begin{enumerate}
  \item The equation \(x^2 + bx + 7 = 0\) has two real solutions. What is the minimum positive integer value of \(b\) for which this is true?
  \item A conical cup has a diameter of 10 cm and a height of 12 cm. What is the volume, in cubic centimeters, of the cup? Express your answer in terms of \(\pi\).
  \item If \(\cos(x^\circ) = \sin(2x + 6)^\circ\), what is the value of \(x\)?
  \item A 15\% acid solution and a 40\% acid solution are mixed to form 50 liters of a 30\% acid solution. How many liters of the 40\% solution are used?
  \item A sphere is inscribed in a cube with edge length 8. What is the volume of the sphere? Express your answer in terms of \(\pi\).
\end{enumerate}

\newpage

\section*{Answer Key and Solutions: Full Mixed Review}

\subsection*{Part A Solutions: Fundamentals Review}
\begin{enumerate}
  \item \(4x - 12 = 2x + 10 \implies 2x = 22 \implies x = 11\).

  \(\boxed{11}\)

  \item \(f(4) = 16 - 12 + 2 = 6\).

  \(\boxed{6}\)

  \item Slope \(= \dfrac{12 - 3}{4 - 1} = \dfrac{9}{3} = 3\).

  \(\boxed{3}\)

  \item \(0.20 \times 150 = 30\).

  \(\boxed{30}\)

  \item \(x + 2x + 60 = 180 \implies 3x = 120 \implies x = 40\).

  \(\boxed{40}\)
\end{enumerate}

\subsection*{Part B Solutions: Core Applications}
\begin{enumerate}
  \item From the second equation, \(x = y + 2\). Substitute: \(3(y + 2) + 2y = 16 \implies 5y + 6 = 16 \implies y = 2\), \(x = 4\).

  \(\boxed{(4, 2)}\)

  \item Sector area \(= \dfrac{45}{360} \cdot \pi(8)^2 = \dfrac{1}{8} \cdot 64\pi = 8\pi\).

  \(\boxed{8\pi}\)

  \item Factor: \(\dfrac{2(x^2 - 4)}{2(x + 2)} = \dfrac{2(x-2)(x+2)}{2(x+2)} = x - 2\).

  \(\boxed{x - 2}\)

  \item Red marbles: \(\dfrac{3}{10} \times 40 = 12\).

  \(\boxed{12}\)

  \item \(c = \sqrt{7^2 + 24^2} = \sqrt{49 + 576} = \sqrt{625} = 25\).

  \(\boxed{25}\)
\end{enumerate}

\subsection*{Part C Solutions: Multi-Step Problems}
\begin{enumerate}
  \item Let the wholesale cost be \(w\). Then \(1.25w = 60 \implies w = 48\).

  \(\boxed{\$48}\)

  \item \(h = -\dfrac{-12}{2(2)} = 3\). \(k = 2(9) - 36 + 22 = 4\). So \(h + k = 7\).

  \(\boxed{7}\)

  \item Shortest side (opposite \(30^\circ\)): \(\dfrac{20}{2} = 10\). Side opposite \(60^\circ\): \(10\sqrt{3}\).

  \(\boxed{10\sqrt{3}}\)

  \item Original sum: \(8 \times 15 = 120\). New sum: \(120 + 33 = 153\). New mean: \(\dfrac{153}{9} = 17\).

  \(\boxed{17}\)

  \item Area ratio \(= \left(\dfrac{2}{3}\right)^2 = \dfrac{4}{9}\). Smaller area: \(54 \times \dfrac{4}{9} = 24\).

  \(\boxed{24}\)
\end{enumerate}

\subsection*{Part D Solutions: Challenging Applications}
\begin{enumerate}
  \item \(\cos(\theta) = \sqrt{1 - 64/289} = \sqrt{225/289} = \dfrac{15}{17}\). Then \(\tan(\theta) = \dfrac{8/17}{15/17} = \dfrac{8}{15}\).

  \(\boxed{\dfrac{8}{15}}\)

  \item Volume of water: \(\pi(4)^2(9) = 144\pi\). In the new container: \(\pi(6)^2 h = 36\pi h = 144\pi \implies h = 4\).

  \(\boxed{4}\)

  \item Set \(-x + c = x^2 + 4x + 5\), giving \(x^2 + 5x + (5 - c) = 0\). For tangency, discriminant \(= 0\):
  \[
  25 - 4(5 - c) = 0 \implies 25 - 20 + 4c = 0 \implies 4c = -5 \implies c = -\frac{5}{4}
  \]

  \(\boxed{-\dfrac{5}{4}}\)

  \item Expression: \(5000(1.06)^t\). After 5 years: \(5000(1.06)^5 = 5000(1.33823) \approx 6691\).

  \(\boxed{6{,}691}\)

  \item The diagonal of the rectangle is the diameter of the circle: \(d = 10\). Let the sides be 6 and \(w\). Then \(6^2 + w^2 = 10^2 \implies w^2 = 64 \implies w = 8\). Area: \(6 \times 8 = 48\).

  \(\boxed{48}\)
\end{enumerate}

\subsection*{Part E Solutions: SAT-Hard}
\begin{enumerate}
  \item For two real solutions, discriminant \(> 0\): \(b^2 - 28 > 0 \implies b^2 > 28 \implies b > \sqrt{28} \approx 5.29\). The minimum positive integer is \(b = 6\).

  (Check: \(36 - 28 = 8 > 0\). Confirmed.)

  \(\boxed{6}\)

  \item Radius \(= 5\). Volume \(= \dfrac{1}{3}\pi(5)^2(12) = 100\pi\).

  \(\boxed{100\pi}\)

  \item Using the complementary relationship: \(\cos(x) = \sin(90 - x)\), so \(90 - x = 2x + 6 \implies 84 = 3x \implies x = 28\).

  \(\boxed{28}\)

  \item Let \(a\) = liters of \(15\%\) and \(b\) = liters of \(40\%\), with \(a + b = 50\) and \(0.15a + 0.40b = 0.30(50) = 15\).
  \[
  0.15(50 - b) + 0.40b = 15 \implies 7.5 + 0.25b = 15 \implies 0.25b = 7.5 \implies b = 30
  \]

  \(\boxed{30}\)

  \item The sphere inscribed in a cube has diameter equal to the cube's edge: \(d = 8\), so \(r = 4\). Volume \(= \dfrac{4}{3}\pi(4)^3 = \dfrac{256\pi}{3}\).

  \(\boxed{\dfrac{256\pi}{3}}\)
\end{enumerate}

\ReflectionPage{Unit 10, Topic 6}{https://www.scotchegg.co}

\end{document}
